\documentclass[12pt]{article}
 
\usepackage[margin=1in]{geometry} 
\usepackage[utf8]{inputenc}
\usepackage[T1]{fontenc}
\usepackage[french]{babel}
\usepackage[autolanguage]{numprint}
\usepackage{lmodern}
\usepackage{amsmath}
\usepackage{amssymb}
\usepackage{cases}
\usepackage{mathrsfs}
\usepackage{textcomp}
\usepackage{stmaryrd}
\usepackage{moreverb}
\usepackage{listings}
\usepackage{color}
\usepackage[usenames,dvipsnames]{xcolor}
\usepackage{graphicx}
\usepackage{fancyvrb}  % Pour le verbatim
\usepackage{fancyhdr}
\pagestyle{fancy}
\usepackage{esint}
\usepackage{tabularx}
\usepackage{multirow}
\usepackage{lipsum}
\usepackage{chngcntr}
\usepackage{changepage}
\usepackage{pdfpages}
\usepackage[shortlabels]{enumitem}
\usepackage{dsfont}
\usepackage{nicefrac}
\usepackage{csquotes}
\usepackage[version=4]{mhchem}
\usepackage{cancel}
\usepackage{eurosym}

\usepackage{siunitx}
\sisetup{
    inter-unit-product = \ensuremath{{}\!\cdot\!{}},
    detect-all,
    separate-uncertainty = true,
    exponent-product = \times,
    space-before-unit = true,
    output-decimal-marker = {,},
    multi-part-units = brackets,
    range-phrase = --,
    % allow-number-unit-breaks,
    list-final-separator = { et },
    list-pair-separator = { et },
    abbreviations
}

\input{../../Template/Commandes.sty}
\begin{document}
 
% --------------------------------------------------------------
%                         Start here
% --------------------------------------------------------------

\lhead{PC Théophile Gautier}
%\chead{L.Villa}
\rhead{Étude documentaire: applications de l'effet tunnel}

\paragraph{Microscopie à effet tunnel}
\begin{enumerate}
\item Citer deux autres exemples en physique où il peut exister une onde évanescente. 
\item Par un calcul d'ordre de grandeur, justifier la phrase du texte: \textit{l'effet tunnel n'est possible que dans le domaine microscopique, pour des particules extrêmement légères et de très faibles barrières de potentiel}.
\item Que représente la couche d'air qui sépare la pointe conductrice de l'échantillon ?
\item Pourquoi est-il important que \textit{la pointe ne comporte que quelques atomes à son extrémité} ? 
\end{enumerate}

\paragraph{Modèle de Gamow de la radioactivité alpha}
\begin{enumerate}
\item En s'aidant du tableau périodique des éléments, écrire la réaction de désintégration $\alpha$ d'un noyau de radium 226. 
\item À partir des masses des noyaux (on rappelle que la masse d'un nucléon vaut environ \SI{1,6e-27}{kg}), en déduire l'énergie libérée par cette transformation nucléaire sous la forme $E=\Delta m \times c^2$ (relation d'Einstein). La calculer d'abord en \SI{}{J} puis en \SI{}{MeV}. Pourquoi cette énergie correspond-elle à l'énergie cinétique de la particule $\alpha$? 
\item On suppose le noyau de radium sphérique, de rayon $R=r_{0}A^{1/3}$ avec $r_0=\SI{1,2e-15}{m}$. En déduire une estimation de $R$ pour le noyau de radium. 
\item À l'aide du tableau ci-dessous, vérifier la loi de Geiger-Nuttal. En gardant les valeurs numériques de $\tau_{1/2}$ en \SI{}{s} et celles de $E$ en \SI{}{MeV}, en déduire une estimation des constantes $A$ et $B$ par une régression linéaire (avec Python ou à la calculatrice).

\begin{center}
\begin{tabular}{c|c|c}
Noyau & $\tau_{1/2}$ (\SI{}{s}) & $E$ (\SI{}{MeV}) \\[1mm]
\hline\hline
\\[-4mm]
$^{212}_{83}\text{Bi}$ & \SI{4,0e4}{} & \SI{6,2}{} \\ [1mm]
$^{212}_{84}\text{Po}$ & \SI{3,0e-7}{} & \SI{9,0}{} \\ [1mm]
$^{215}_{85}\text{At}$ & \SI{1,0e-4}{} & \SI{8,1}{} \\ [1mm]
$^{222}_{86}\text{Ra}$ & \SI{3,3e5}{} & \SI{5,6}{} \\ 
\end{tabular} 
\qquad
\begin{tabular}{c|c|c}
Noyau & $\tau_{1/2}$ (\SI{}{s}) & $E$ (\SI{}{MeV}) \\[1mm]
\hline\hline
\\[-4mm]
$^{226}_{88}\text{Ra}$ & \SI{5,4e10}{} & \SI{4,9}{} \\ [1mm]
$^{232}_{90}\text{Th}$ & \SI{4,4e17}{} & \SI{4,0}{} \\ [1mm]
$^{236}_{92}\text{U}$ & \SI{7,2e14}{} & \SI{4,4}{} \\ [1mm]
$^{242}_{96}\text{Cm}$ & \SI{1,4e7}{} & \SI{6,2}{} \\ 
\end{tabular} 
\end{center}
\item Estimer la demi-vie du radium 226 prédite par le modèle de Gamow et comparer à la valeur expérimentale figurant dans le tableau ci-dessus. Commenter. 
\end{enumerate}

\end{document}